%══════════════════════════════════════════════════════════════════════
%  CHAPTER 4 — Proposed System: PAD-ONAP
%══════════════════════════════════════════════════════════════════════
\chapter{Proposed System: PAD-ONAP}
\label{ch:system}

This chapter presents the \padonap{} system in full. It begins with the
design rationale (Section~\ref{sec:approach}), proceeds through the
four-stage pipeline architecture (Section~\ref{sec:arch-overview}),
formalises the threat model and safety invariants
(Sections~\ref{sec:threat-model}--\ref{sec:safety}), and details the
five-tier policy state machine, LP SLA allocator, and AIOutputPayload
interface contract (Sections~\ref{sec:tiers}--\ref{sec:payload}).

%%─────────────────────────────────────────────────────────────────────
\section{Approach Direction}
\label{sec:approach}

The central design insight motivating \padonap{} is as follows.
Any reactive DDoS mitigation system has a hard latency floor equal to
the VNF boot time: the orchestrator cannot act faster than its slowest
required function. For a traffic scrubber this floor is approximately
6~s (Section~\ref{sec:nfv-bg}). No amount of orchestration tuning lowers
this floor.

The only way to reduce \emph{effective} mitigation latency is to
pre-position the mitigation VNF \emph{before} the attack reaches the
reactive trigger. This requires a forecast. The forecast trades reduced
activation latency (VNF already running when attack peaks) against a
modest false-positive risk (the VNF consumes resources on benign windows
when the forecast was incorrect).

The \padonap{} design addresses this trade-off with three architectural
decisions:

\begin{enumerate}
  \item \textbf{Two-tier proactive action.} Pre-positioning uses
        \tier{2} (rate-limiter, 505~ms boot) rather than \tier{3}
        (scrubber, 6{,}006~ms boot). If the forecast is wrong, the
        cost is 500~ms wasted boot time and moderate resource consumption.
        If correct, the rate-limiter is already in place when the attack
        peaks.

  \item \textbf{LP SLA allocator at every tier transition.} Activating
        any VNF subtracts overhead from the available link capacity.
        An explicit LP solve at each tier change ensures that all tenant
        SLA floors are respected regardless of which tier is active.

  \item \textbf{Hysteresis guard.} A hold window of $k{=}2$ prevents
        rapid oscillation between tiers (VNF churn), which would add
        latency variance and increase resource overhead.
\end{enumerate}

%%─────────────────────────────────────────────────────────────────────
\section{System Overview}
\label{sec:arch-overview}

\padonap{} is a four-stage pipeline. Every 5-second window, stage~S1
scrapes per-device gNMI counters and (optionally) NetFlow v5 records;
stage~S2 aggregates them into a 17-dimensional feature vector
(Section~\ref{sec:features}); stage~S3 runs XGBoost 7-class detection,
Transformer+LSTM 4-horizon forecasting, and emits a canonical
\texttt{AIOutputPayload} JSON; stage~S4 consumes this payload, evaluates
the five-tier policy, allocates tenant bandwidth under an LP SLA model,
and instantiates VNFs through an ONAP SO client.

\begin{figure}[H]
\centering
\begin{tikzpicture}[
  stage/.style={draw, rounded corners, minimum width=2.8cm,
                minimum height=1.1cm, align=center, font=\small},
  arr/.style={-Latex, thick},
]
\node[stage, fill=blue!10]  (s1) {S1\\Telemetry\\(gNMI, NetFlow)};
\node[stage, fill=green!10, right=0.55cm of s1] (s2)
      {S2\\Features\\(17 entropy/rate)};
\node[stage, fill=orange!15, right=0.55cm of s2] (s3)
      {S3 AI\\(XGB +\\Txf-LSTM)};
\node[stage, fill=red!12, right=0.55cm of s3] (s4)
      {S4 Orch.\\(Policy+SLA\\+ONAP SO)};
\draw[arr] (s1)--(s2);
\draw[arr] (s2)--(s3);
\draw[arr] (s3)--(s4);
\node[below=0.6cm of s4, text width=3.5cm, align=center,
      font=\small\itshape]
  {CLAMP/DMaaP\\feedback to S3};
\draw[arr, dashed, bend right=30] (s4.south) to
  node[below, font=\tiny] {} (s3.south);
\end{tikzpicture}
\caption{\padonap{} four-stage pipeline. Each stage is independently
replaceable via the \texttt{AIOutputPayload} JSON contract on DMaaP.}
\label{fig:arch}
\end{figure}

The DMaaP message bus between S3 and S4 (and the feedback from S4 back
to S3 via CLAMP) is the key architectural boundary. Because all inter-stage
communication passes through the canonical JSON schema on a named Kafka
topic, any stage can be upgraded or replaced independently. This is the
``separation of concerns'' property claimed in Contribution~1.

%%─────────────────────────────────────────────────────────────────────
\section{Threat Model}
\label{sec:threat-model}

We formalise the adversary capability and the system's trust assumptions
before stating the safety invariants.

\begin{assumption}[Telemetry integrity]
\label{asm:telemetry}
The gNMI/NetFlow path from network devices to stage~S2 is authenticated
and integrity-protected (e.g.~gNMI over mutual-TLS). The adversary cannot
forge or modify telemetry in transit. Compromise of the telemetry
infrastructure itself is out of scope.
\end{assumption}

\begin{assumption}[Control-plane trust]
\label{asm:control}
The ONAP SO/CLAMP/Policy plane and the DMaaP bus are operated within a
trusted administrative domain. Compromise of the orchestrator control
plane is out of scope.
\end{assumption}

\begin{definition}[Adversary capability]
\label{def:adv}
The adversary $\mathcal{A}$ controls a bounded-size botnet and may:
(a) emit volumetric floods (UDP, ICMP, amplification) up to link
saturation; (b) execute protocol-exhaustion attacks (SYN flood);
(c) ramp attack rate gradually to evade static rate thresholds; and
(d) sustain low-and-slow attacks that remain below a specified rate.
$\mathcal{A}$ cannot forge gNMI counters, inject messages on DMaaP, or
observe the ML model parameters.
\end{definition}

Out-of-scope threats include insider attacks, cryptographic attacks on
TLS, and poisoning of the AI training data. The XGBoost model is trained
offline on CICDDoS2019; data-poisoning defences are left for future
work (Section~\ref{sec:future}).

%%─────────────────────────────────────────────────────────────────────
\section{Safety Definition}
\label{sec:safety}

Proactive mitigation changes the relationship between the defender and
the tenant: a tier may be raised \emph{before} any attack has actually
crossed a detection threshold. Three safety properties are required.

\begin{definition}[Pre-positioning safety]
\label{def:safety}
For every tier transition $T_i \to T_j$ triggered proactively (i.e.~driven
by a forecast output, not by a detection-confidence crossing), the
bandwidth allocation to every benign tenant $u$ must satisfy
\[
  a_u(T_j) \;\geq\; \text{floor}_u,
\]
where $\text{floor}_u$ is the SLA floor of tenant $u$ and $a_u(T_j)$
is its LP-allocated bandwidth at tier $T_j$.
\end{definition}

\begin{invariant}[URLLC floor]
\label{inv:urllc}
For every mitigation tier $T \in \{T_0, T_1, T_2, T_3, T_4\}$ and in
every evaluated scenario,
\[
  a_{\text{URLLC}}(T) \;\geq\; 150\;\text{Mbps}.
\]
\end{invariant}

\begin{invariant}[Bounded rollback]
\label{inv:rollback}
After a proactive escalation $T_i \to T_j$ at window $w$, if the raw
tier output reports $T_k < T_j$ for $k{=}2$ consecutive windows, the
policy engine de-escalates within window $w + 1 + k$. No proactive
escalation is held longer than $k$ windows beyond the last supporting
signal.
\end{invariant}

Both Invariants~\ref{inv:urllc} and~\ref{inv:rollback} are verified
empirically in Chapter~\ref{ch:eval} across all eight evaluation
scenarios.

%%─────────────────────────────────────────────────────────────────────
\section{Five-Tier Policy State Machine}
\label{sec:tiers}

The mitigation action space is discretised into five tiers
(Table~\ref{tab:tiers}). Tier selection is a deterministic function of
the AI payload's top-horizon forecast probability and detection
classification confidence, passed through a hysteresis filter and a
frequency guard.

\begin{table}[H]
\centering
\small
\caption{Five-tier mitigation ladder. VNF profiles are defined in
\texttt{vnf\_catalogue/}.}
\label{tab:tiers}
\begin{tabular}{@{}clll@{}}
\toprule
Tier & Label & Action & VNF profile \\
\midrule
$T_0$ & NORMAL  & No action                           & --- \\
$T_1$ & ALERT   & Log + raise operator alert           & --- \\
$T_2$ & PREEMPT & Pre-position rate-limiter VNF        & vnfd-ratelimiter-v1 \\
$T_3$ & MITIGATE& Instantiate scrubber VNF + SFC chain & vnfd-scrubber-v1 \\
$T_4$ & SEVERE  & Blackhole upstream + scrubber        & vnfd-blackhole-v1 \\
\bottomrule
\end{tabular}
\end{table}

\paragraph{Hysteresis filter.}
Let $\tau(w)$ denote the raw tier emitted by the AI engine at window $w$.
The published tier $\tilde{\tau}(w)$ is computed as
\[
  \tilde{\tau}(w) \;=\;
  \max\!\Bigl(\tau(w),\;
              \tilde{\tau}(w{-}1)
              \cdot \mathbf{1}\bigl[w - w_{\mathrm{last}} \leq k\bigr]\Bigr),
\]
where $k{=}2$ is the hold window and $w_{\mathrm{last}}$ is the last
window at which the raw tier equalled $\tilde{\tau}(w{-}1)$. This
guarantees that the published tier cannot drop by more than one level
per window, implementing Invariant~\ref{inv:rollback}.

\paragraph{Frequency guard.}
To prevent rapid VNF churn on noisy signals, the policy engine records
the time of the last instantiation event. A new instantiation at the same
tier is blocked if the elapsed time is less than a configurable guard
interval (default: 10 windows = 50~s). This guard is \emph{disabled}
in evaluation mode (\texttt{eval\_mode=True}) to allow scenario replay.

%%─────────────────────────────────────────────────────────────────────
\section{LP Multi-Tenant SLA Allocator}
\label{sec:sla-lp}

Let $n$ tenants share a total link capacity $C$ (Mbps). At tier $T$
the VNF chain consumes an overhead $o(T)$ Mbps, leaving $C - o(T)$
for tenants. Denote by $d_u$ tenant $u$'s demand, $f_u$ its guaranteed
floor, and $w_u$ its priority weight. The SLA allocator solves the LP

\begin{align}
  \max_{\mathbf{a}} &\quad \sum_{u=1}^{n} w_u \, a_u
  \label{eq:lp-obj} \\
  \text{s.t.} &\quad \sum_{u=1}^{n} a_u \;\leq\; C - o(T),
  \label{eq:lp-cap} \\
  &\quad f_u \;\leq\; a_u \;\leq\; d_u,
  \qquad \forall\, u \in \{1,\ldots,n\}.
  \label{eq:lp-bounds}
\end{align}

The objective~\eqref{eq:lp-obj} maximises weighted total allocation,
corresponding to best-effort service for higher-priority tenants.
Constraint~\eqref{eq:lp-cap} enforces the total capacity budget.
Constraint~\eqref{eq:lp-bounds} guarantees per-tenant floor and cap;
for $u = \text{URLLC}$, $f_u = 150$~Mbps directly encodes
Invariant~\ref{inv:urllc}.

The LP is solved with HiGHS through SciPy
(\texttt{scipy.optimize.linprog}). If SciPy is unavailable, a
proportional fallback allocates capacity in proportion to demand
weights, subject to floors. In the three-tenant configuration
evaluated in Scenario~S7: $n = 3$, $C = 1000$~Mbps,
$f_{\text{URLLC}} = 150$, $f_{\text{eMBB}} = 100$,
$f_{\text{mMTC}} = 50$; $d_u = 500$ for all tenants; overhead at
$T_2$ is 500~Mbps.

%%─────────────────────────────────────────────────────────────────────
\section{AIOutputPayload Interface Contract}
\label{sec:payload}

Stage~S3 emits a JSON object on the DMaaP topic
\texttt{PAD\_ONAP\_AI\_SIGNALS} at the end of every window. The schema
is versioned and fully specified in
\texttt{pipeline/s3\_ai/payload\_schema.json}. The skeleton is:

\begin{lstlisting}[language=json, caption={Skeleton of AIOutputPayload
v1. All numeric fields are rounded to 4 decimal places.}]
{
  "schema_version": "1.0",
  "timestamp":      <epoch seconds, float>,
  "device_id":      "<string>",
  "detection": {
    "attack_class": <int, 0..6>,
    "confidence":   <float, 0..1>
  },
  "forecast": [
    {"horizon_s": 30,  "p_attack": <float>},
    {"horizon_s": 60,  "p_attack": <float>},
    {"horizon_s": 90,  "p_attack": <float>},
    {"horizon_s": 120, "p_attack": <float>}
  ],
  "top_features": [
    {"feature": "<name>", "shap": <float>},
    ...
  ],
  "suggested_tier": <int, 0..4>,
  "proactive":      <bool>
}
\end{lstlisting}

The \texttt{proactive} flag is \texttt{true} when the tier suggestion
originates from the forecaster path rather than the detector path. Stage~S4
logs this flag per window, enabling the lead-time analysis described in
Chapter~\ref{ch:eval}. Downstream ONAP consumers need parse only this
schema; the AI implementation can be replaced without any change to
the policy engine.

%%─────────────────────────────────────────────────────────────────────
\section{Chapter Summary}
\label{sec:system-summary}

This chapter presented the full \padonap{} system design. The four-stage
pipeline (S1 telemetry $\to$ S2 features $\to$ S3 AI $\to$ S4 orchestration)
is connected by a canonical \texttt{AIOutputPayload} JSON schema on DMaaP.
The threat model bounds the adversary to botnet-level volumetric and
protocol attacks. Three formal safety properties---pre-positioning safety,
URLLC floor invariant, and bounded rollback---define the correctness
conditions that the system must satisfy. The five-tier policy state machine
with hysteresis and the LP SLA allocator operationalise these conditions
at runtime. The AI components that drive stage~S3 are detailed in the
next chapter.
